\documentclass[11pt,a4paper]{article}
\usepackage[utf8]{inputenc}
\usepackage[T1]{fontenc}
\usepackage[spanish]{babel}
\usepackage{geometry}
\geometry{margin=2.5cm}

\title{Las tres escuelas del átomo}
\author{}
\date{}

\begin{document}
\maketitle

\section*{Quiénes son}

Esta es la historia de cómo tres tradiciones científicas distintas ---la física nuclear experimental británica, la física teórica alemana y la física de partículas italiana--- convergieron para crear la imagen moderna del átomo. Tres hombres, cada uno en su país y su estilo, fundaron escuelas que formaron a la siguiente generación. Sus alumnos se cruzaron, competieron y, a veces, colaboraron.

Ernest Rutherford nació en 1871 en Brightwater, Nueva Zelanda, hijo de un molinero escocés y una maestra inglesa. Estudió en el Canterbury College de Christchurch, donde destacó tanto en ciencias como en rugby. Era corpulento, de voz potente y risa fácil; se dice que su carcajada se escuchaba en los pasillos del Cavendish Laboratory de Cambridge, donde llegó en 1895 gracias a una beca. Bajo la dirección de J. J. Thomson, descubrió que la radiactividad implicaba la transformación de un elemento en otro. En 1907 se trasladó a la Universidad de Manchester, donde construyó un laboratorio que rivalizaba con Cambridge. Le gustaba el golf, aunque jugaba mal, y tenía fama de despreciar la física teórica abstracta: solía decir que los modelos matemáticos eran útiles solo si podías ponerles un cable y hacer que funcionaran.

Niels Bohr nació en Copenhague en 1885. Su padre era fisiologista, su hermano mayor matemático. Bohr era lento redactando, torpe en los exámenes formales y famoso por necesitar que otros ---especialmente matemáticos--- le pulieran los cálculos. Estudió en Copenhague y en Cambridge, donde en 1912 conoció a Rutherford. De esa estancia nació el modelo atómico de Bohr, publicado en 1913, donde los electrones giran en órbitas cuantizadas. En 1921 fundó el Instituto de Física Teórica de Copenhague, un edificio modesto con calefacción irregular pero una atmósfera intelectual única. Bohr era un conversador obstinado: podía discutir una idea durante horas, reformulándola hasta que el interlocutor ya no sabía si estaba de acuerdo o en desacuerdo.

Werner Heisenberg nació en Wurzburgo en 1901. Su padre era profesor de griego clásico en la Universidad de Múnich, y Heisenberg creció leyendo Platón en el original. Estudió física en Múnich con Arnold Sommerfeld, pero en 1923 se trasladó a Gotinga para trabajar con Max Born. Era un pianista competente, casi profesional, y durante toda su vida mantuvo que la física y la música eran dos formas del mismo impulso: encontrar la estructura oculta. En 1924 viajó a Copenhague con una beca; Bohr lo adoptó como discípulo filosófico. En 1925, en la isla de Heligoland, desarrolló la mecánica de matrices. Tenía 24 años.

Lev Landau nació en Bakú en 1908, en el seno de una familia judía de ingenieros petroleros. Era un niño prodigio que se graduó a los 14 años y se doctoró a los 19 en la Universidad de Leningrado. En 1929 viajó a Europa Occidental: estuvo en Copenhague con Bohr, en Cambridge con Dirac, y en 1930 en Leipzig con Heisenberg. Landau era famoso por su examen oral destructivo, conocido como el mínimo de Landau: un test de conocimientos que solo un puñado de físicos soviéticos lograron pasar. Era también un coleccionista de anécdotas, clasificadas por números, y un ateo militante que en 1938 fue encarcelado durante un año por criticar al régimen estalinista.

Max Born nació en Breslavia (hoy Wrocław, Polonia) en 1882. Su padre era médico, su madre murió cuando él tenía cuatro años. Estudió en Gotinga, donde se doctoró con Carl Runge, y luego en Cambridge con J. J. Thomson. En 1921 fue nombrado director del Instituto de Física Teórica de la Universidad de Gotinga, que convirtió en el centro mundial de la mecánica cuántica. Born era un hombre melancólico, propenso a la depresión, que encontraba consuelo en la música de Bach y en largas caminatas por el Harz. En 1954 recibió el Nobel de Física por su interpretación probabilística de la función de onda, aunque siempre lamentó que la prioridad del descubrimiento se le discutiera.

Wolfgang Pauli nació en Viena en 1900. Su padre, Wolfgang Joseph Pauli, era químico y profesor en la Universidad de Viena; su madre, Bertha Schütz, era periodista y corresponsal del diario Neue Freie Presse. Pauli era un niño precoz que publicó artículos sobre relatividad general a los 18 años. Estudió en Múnich con Sommerfeld, pero en 1921 se trasladó a Gotinga como asistente de Max Born. Era famoso por su crítica mordaz: decían que su sola presencia en un laboratorio hacía que los aparatos dejaran de funcionar. Creía en la sincronicidad y mantenía una correspondencia larga con Carl Jung sobre el inconsciente colectivo. En 1945 recibió el Nobel por el principio de exclusión.

J. Robert Oppenheimer nació en Nueva York en 1904, en una familia judía acomodada. Su padre importaba telas, su madre era pintora. Estudió en la Ethical Culture School, en Harvard (donde cursó química, física, filosofía y griego clásico simultáneamente), y en 1925 viajó a Europa. En Cambridge intentó trabajar con J. J. Thomson, pero fue un desastre: una depresión lo llevó a una crisis nerviosa. Se recuperó en Gotinga, donde Max Born lo aceptó como estudiante de doctorado. Oppenheimer leía sánscrito fluídamente y recitaba pasajes del Bhagavad Gita mientras conducía. En 1943 fue nombrado director del Laboratorio de Los Álamos.

Enrico Fermi nació en Roma en 1901. Su padre era funcionario de ferrocarriles, su madre maestra de escuela. Era un niño autodidacta que aprendió física de libros de texto encontrados en un mercado de pulgas. En 1922 se doctoró en Pisa con una tesis sobre rayos X. De regreso en Roma, en 1927 fue nombrado profesor de física teórica en la Universidad de Roma, donde reunió a un grupo de jóvenes físicos conocidos como los Ragazzi di Via Panisperna, por la calle donde estaba el Instituto de Física. Fermi era metódico, pausado, de apariencia modesta. Jugaba al tenis los domingos y tenía una memoria fotográfica para los números.

Emilio Segrè nació en Tívoli, cerca de Roma, en 1905. Estudió ingeniería en la Universidad de Roma La Sapienza, donde en 1927 conoció a Fermi. Se unió al grupo de Via Panisperna como el más joven. Era de familia judía sefardí; en 1938, tras las leyes raciales de Mussolini, fue despedido de la universidad y emigró a Berkeley, donde descubrió el antiprotón en 1955 (Nobel 1959). Era un hombre de temperamento áspero, famoso por sus disputas académicas.

Edoardo Amaldi nació en Carpaneto Piacentino en 1908. Estudió en Roma, donde en 1928 conoció a Fermi y se unió al grupo de Via Panisperna. Era el más equilibrado del grupo: mientras Fermi era reservado y Segrè irascible, Amaldi era conciliador y organizador. Sobrevivió a la guerra en Roma, escondiendo los aparatos de radiactividad del régimen fascista. Después de 1945, fue el principal artífice de la reconstrucción de la física italiana y uno de los fundadores del CERN.

\section*{La historia}

1912. Niels Bohr llega a Manchester con una beca del Carlsberg Foundation. Ernest Rutherford lo recibe en su laboratorio, le muestra los experimentos de dispersión de partículas alfa y lo desafía a aplicar las ideas cuánticas de Planck al átomo. Bohr trabaja día y noche en una habitación alquilada cerca del laboratorio. En julio de 1912 envía a Rutherford un borrador de lo que sería su modelo atómico. Rutherford lo corrige, le sugiere omitir una sección sobre la química, y lo anima a publicar. El artículo aparece en 1913. Bohr siempre consideró a Rutherford su padre científico.

1924. Werner Heisenberg, entonces estudiante de 23 años, recibe una beca de la Rockefeller Foundation para trabajar en Copenhague. Niels Bohr lo espera en la estación; van directamente al Instituto y comienzan a discutir. Durante los siguientes meses, Bohr y Heisenberg pasean por los jardines de Frederiksberg hablando de física y filosofía. Bohr le enseña a Heisenberg que la mecánica cuántica no es solo matemática: es una nueva forma de preguntarse qué significa medir. Heisenberg, a su vez, le enseña a Bohr las técnicas matemáticas de Gotinga. En 1925, Heisenberg desarrolla la mecánica de matrices; Bohr lo revisa y lo aprueba.

1930. Lev Landau llega a Leipzig con una beca soviética. Werner Heisenberg lo recibe en su instituto y le asigna un problema sobre el diamagnetismo de los metales. Landau resuelve el problema en semanas, pero su estilo abrasivo genera fricción: corrige a Heisenberg en seminarios públicos y se niega a usar la notación que Heisenberg prefiere. A pesar de la tensión, Heisenberg lo guía en la redacción del artículo que se publicará en Zeitschrift für Physik. Landau siempre citó a Heisenberg como quien le enseñó a hacer física de verdad, no solo matemáticas.

1921--1922. Wolfgang Pauli, recién graduado en Múnich con una tesis sobre la relatividad general que Einstein elogió, llega a Gotinga como asistente de Max Born. Born le encarga revisar los cálculos de un artículo sobre cristales y descubre que Pauli es más rápido y riguroso que él mismo. Durante un año, Born le enseña a Pauli la física del estado sólido y la teoría de perturbaciones. Pauli, a su vez, critica sin piedad los errores de Born. La relación es intensa y fructífera: Born siempre dijo que Pauli fue su conciencia matemática.

1923--1924. Werner Heisenberg se traslada de Múnich a Gotinga para completar su habilitación bajo la dirección de Max Born. Born lo introduce en la teoría de perturbaciones y los métodos matriciales que serán clave para la mecánica cuántica. Heisenberg asiste a los seminarios de Born, donde conoce a Pauli, Jordan y Hund. En 1924, antes de ir a Copenhague, Heisenberg presenta su habilitación sobre el efecto Zeeman anómalo, supervisada por Born. Aunque Heisenberg luego trabajará intensamente con Bohr, Born fue su director formal en Gotinga y coautor de los artículos fundacionales de 1925.

1926--1927. J. Robert Oppenheimer llega a Gotinga tras su fracaso en Cambridge. Max Born lo acepta como estudiante de doctorado, aunque al principio desconfía: Oppenheimer es arrogante, interrumpe a los profesores y habla más de lo que escucha. Born lo pone a trabajar en la teoría cuántica de moléculas. Oppenheimer aprende a calcular, a callar y a escribir. La tesis, sobre la aproximación de Born-Oppenheimer para separar movimientos nucleares y electrónicos, se convierte en un clásico. Oppenheimer siempre citó a Born como quien le enseñó a ser un físico en lugar de un intelectual disperso.

1928--1930. Enrico Fermi, ya profesor en Roma, selecciona a los mejores estudiantes de física de La Sapienza. Emilio Segrè, entonces estudiante de ingeniería, asiste a una conferencia de Fermi sobre la estadística de Fermi-Dirac y queda fascinado. Pide cambiarse a física y Fermi lo acepta en el grupo de Via Panisperna. Segrè trabaja con Fermi en los experimentos de radiactividad inducida con neutrones, que llevarán al descubrimiento de los elementos transuránicos. Fermi supervisa personalmente los cálculos de Segrè y le enseña a diseñar experimentos.

1928--1930. Edoardo Amaldi conoce a Fermi en 1928, en su segundo año de universidad. Fermi lo invita al grupo de Via Panisperna, donde Amaldi se convierte en el organizador: registra los datos, coordina los turnos de los experimentos con neutrones y redacta los artículos. Fermi le enseña a Amaldi la física experimental y la importancia de la precisión numérica. Amaldi siempre dijo que Fermi le transmitió el gusto por lo concreto, por lo que se puede medir.

\end{document}
